\section{Introduction}
\label{sec:intro}

Knowledge graphs (KGs) encode facts as structured triples $(h, r, t)$ and have become foundational infrastructure for knowledge-based systems, powering applications from question answering and recommendation to clinical decision support and drug repurposing~\cite{shen2022comprehensive,biswas2023knowledge}.
When deployed in safety-critical domains---pharmacovigilance, clinical decision support, financial compliance---these systems must distinguish between knowledge that can be acted upon and statistical guesses that cannot.
Yet the tools currently available for evaluating KG completion models provide no mechanism for making this distinction.

\emph{Link prediction}---inferring missing facts from an existing graph---is the central computational task in KG completion.
A large family of embedding models has been developed for this task~\cite{bordes2013translating,sun2019rotate,trouillon2016complex,dettmers2018convolutional,balazevic2019tucker,nguyen2018convkb,sun2022hake,le2023knowledge}, and their evaluation has standardized around two aggregate metrics: Mean Reciprocal Rank (MRR) and Hits@$K$.
These metrics measure how highly the correct tail entity is ranked among all candidates, condensing thousands of test triples into a single scalar.
This protocol has enabled reproducible benchmarking and catalyzed rapid progress over the past decade~\cite{wang2017knowledge,rossi2021knowledge,li2023knowledge,liu2022understanding}.

But a single scalar creates a dangerous blind spot for knowledge-based systems.
A model reporting MRR $= 0.56$ says nothing about \emph{which} of its predictions can be trusted, \emph{which relations} it handles reliably, or \emph{what fraction} of its output is noise.
A pharmaceutical researcher using KG completion to prioritize candidate drug--target interactions cannot determine whether the top-10 candidates are $90\%$ reliable or $10\%$ reliable.
This distinction directly determines experimental resource allocation---and it is invisible to current evaluation protocols.
Recent surveys have identified this as a critical open problem, calling for statistical rigor and reliability quantification in KG evaluation~\cite{jarnac2025uncertainty,biswas2023knowledge}.

We introduce a statistical quality assurance framework that treats KG completion evaluation as a large-scale multiple hypothesis testing problem.
Each test triple corresponds to a null hypothesis $H_0$ asserting that the triple is absent from the true KG; the embedding model assigns each candidate a score that serves as a test statistic.
From this formulation, we construct empirical $p$-values via permutation-based negative sampling---a procedure that requires no distributional assumptions and no model modification---and apply three complementary FDR control procedures: Benjamini-Hochberg FDR control~\cite{benjamini1995controlling}, Storey's $q$-value estimation with $\hat{\pi}_0$ estimation~\cite{storey2002direct}, and Efron's empirical Bayes local FDR~\cite{efron2004large}.
These well-established statistical tools together produce multi-resolution reliability profiles: global diagnostics (what fraction of predictions carry genuine signal?), per-relation stratification (which relations are reliable?), and per-triple calibration (which specific predictions are trustworthy?).
The framework operates as a post-hoc layer on any trained KG embedding model, requiring only the scoring function.

This perspective is grounded in the knowledge-driven methodology central to the KBS community.
Recent KBS publications have advanced KG completion through architectural innovation~\cite{fang2025knowledge,shang2024learnable,lu2024deep}, embedding quality assessment~\cite{purohit2025vanilla}, and knowledge representation~\cite{le2023knowledge,zhang2023conflict}.
Our work extends this line by introducing statistical quality assurance for KG inferences---a capability directly relevant to the reliability needs of deployed knowledge-based systems.

We make the following contributions:

\begin{enumerate}
    \item We formalize KG completion evaluation as a multiple hypothesis testing problem, providing a statistical foundation for per-prediction reliability assessment (Section~\ref{sec:method}).

    \item We develop an empirical $p$-value construction procedure that converts uncalibrated embedding scores into valid $p$-values via permutation-based negative sampling, imposing no distributional assumptions (Section~\ref{subsec:pvalue}).

    \item We introduce a three-level FDR-based evaluation protocol---global diagnostics ($\hat{\pi}_0$), relation-stratified analysis, and cross-model reliability comparison---that reveals dimensions of model quality inaccessible to aggregate ranking metrics (Section~\ref{subsec:protocol}).

    \item We evaluate the framework on WN18RR and FB15k-237 across four model architectures (TransE, RotatE, ComplEx, ConvE). The experiments yield three findings: (i) $\hat{\pi}_0 = 0.259$--$0.851$ across models and benchmarks, indicating 26--85\% of test triples are statistically indistinguishable from noise; (ii) per-relation discovery rates span from below $10\%$ to above $95\%$ within a single model, exposing extreme heterogeneity in prediction reliability; and (iii) model rankings by MRR and by FDR power capture distinct dimensions of quality---RotatE leads on MRR while TransE leads on statistically significant discoveries, and neither metric alone provides a complete reliability profile for model deployment (Section~\ref{sec:experiments}).
\end{enumerate}
